A percepção dos lugares e a construção de diagnósticos sobre o território são aspectos cada vez mais mediados por tecnologias digitais e por dados produzidos espontaneamente pelos cidadãos. Esta dissertação investiga como a análise de sentimento avançada, apoiada por Modelos de Linguagem de Grande Porte (LLMs), pode aprimorar a interpretação dessas experiências subjetivas e transformá-las em informações úteis para o acompanhamento de Objetivos de Desenvolvimento Sustentável (ODS) em escala municipal. A pesquisa adota a metodologia da Design Science para orientar a concepção, implementação e avaliação de um artefato computacional capaz de coletar avaliações públicas georreferenciadas, aplicar filtros de relevância geográfica e semântica, analisar sentimentos por aspecto e gerar diagnósticos estruturados em linguagem natural. A abordagem utiliza dados da Google Places API e modelos como Gemini e Perplexity para analisar percepções associadas a serviços, instituições, infraestruturas urbanas e demais locais relacionados aos ODS. O objetivo principal é responder à questão: como a análise de sentimento por aspecto, aplicada a dados perceptivos georreferenciados, pode auxiliar a identificar gargalos no desempenho de ODS municipais e subsidiar políticas públicas? A aplicação no contexto de Mossoró-RN demonstra o potencial da metodologia para revelar padrões recorrentes em comentários públicos, identificar dimensões críticas em diferentes ODS, apoiar a construção de dashboards interativos e oferecer subsídios para gestores públicos sem exigir expertise técnica avançada. Assim, a pesquisa busca validar um método escalável para capturar, interpretar e transformar percepções cidadãs sobre lugares em diagnósticos acionáveis para o planejamento urbano, a gestão pública e o desenvolvimento sustentável.

% Separe as palavras-chave por ponto
\palavraschave{Objetivos de Desenvolvimento Sustentável. Percepção de Lugar. Processamento de Linguagem Natural. Análise de Sentimentos. Modelos de Linguagem de Grande Porte. Google Places API. Políticas Públicas.}
